\documentclass[12pt,a4paper]{article}

\usepackage[T2A]{fontenc}
\usepackage[utf8]{inputenc}
\usepackage[russian]{babel}
\usepackage{amsmath,amssymb,amsthm,mathtools}
\usepackage{enumitem}
\usepackage{listings}
\usepackage{xcolor}
\usepackage{geometry}

\geometry{margin=2.2cm}

\newtheorem{corollary}{Следствие}
\newtheorem{definition}{Определение}
\newtheorem{theorem}{Теорема}
\newtheorem{lemma}{Лемма}
\newtheorem{example}{Пример}
\newtheorem{remark}{Замечание}

\lstdefinestyle{fpstyle}{
  basicstyle=\ttfamily\small,
  keywordstyle=\color{blue},
  commentstyle=\color{gray},
  stringstyle=\color{teal},
  showstringspaces=false,
  frame=single,
  breaklines=true,
  tabsize=2,
  columns=fullflexible
  literate=
    {А}{{\CYRA}}1
    {Б}{{\CYRB}}1
    {В}{{\CYRV}}1
    {Г}{{\CYRG}}1
    {Д}{{\CYRD}}1
    {Е}{{\CYRE}}1
    {Ё}{{\CYRYO}}1
    {Ж}{{\CYRZH}}1
    {З}{{\CYRZ}}1
    {И}{{\CYRI}}1
    {Й}{{\CYRISHRT}}1
    {К}{{\CYRK}}1
    {Л}{{\CYRL}}1
    {М}{{\CYRM}}1
    {Н}{{\CYRN}}1
    {О}{{\CYRO}}1
    {П}{{\CYRP}}1
    {Р}{{\CYRR}}1
    {С}{{\CYRS}}1
    {Т}{{\CYRT}}1
    {У}{{\CYRU}}1
    {Ф}{{\CYRF}}1
    {Х}{{\CYRH}}1
    {Ц}{{\CYRC}}1
    {Ч}{{\CYRCH}}1
    {Ш}{{\CYRSH}}1
    {Щ}{{\CYRSHCH}}1
    {Ъ}{{\CYRHRDSN}}1
    {Ы}{{\CYRERY}}1
    {Ь}{{\CYRSFTSN}}1
    {Э}{{\CYREREV}}1
    {Ю}{{\CYRYU}}1
    {Я}{{\CYRYA}}1
    {а}{{\cyra}}1
    {б}{{\cyrb}}1
    {в}{{\cyrv}}1
    {г}{{\cyrg}}1
    {д}{{\cyrd}}1
    {е}{{\cyre}}1
    {ё}{{\cyryo}}1
    {ж}{{\cyrzh}}1
    {з}{{\cyrz}}1
    {и}{{\cyri}}1
    {й}{{\cyrishrt}}1
    {к}{{\cyrk}}1
    {л}{{\cyrl}}1
    {м}{{\cyrm}}1
    {н}{{\cyrn}}1
    {о}{{\cyro}}1
    {п}{{\cyrp}}1
    {р}{{\cyrr}}1
    {с}{{\cyrs}}1
    {т}{{\cyrt}}1
    {у}{{\cyru}}1
    {ф}{{\cyrf}}1
    {х}{{\cyrh}}1
    {ц}{{\cyrc}}1
    {ч}{{\cyrch}}1
    {ш}{{\cyrsh}}1
    {щ}{{\cyrshch}}1
    {ъ}{{\cyrhrdsn}}1
    {ы}{{\cyrery}}1
    {ь}{{\cyrsftsn}}1
    {э}{{\cyrerev}}1
    {ю}{{\cyryu}}1
    {я}{{\cyrya}}1
}

\lstset{style=fpstyle}

\title{Билет №47 \\ \small Особенности программирования систем реального времени. (СпБПУ)}
\author{Сериков Владислав Александрович}
\date{13 мая 2026}

\begin{document}

\maketitle
\tableofcontents
\newpage

\section{Особенности программирования систем реального времени}

\subsection{Понятие системы реального времени}

\textbf{Система реального времени} --- это вычислительная система, корректность которой определяется не только логической правильностью результата, но и временем, к которому этот результат должен быть получен.

Иначе говоря, программа реального времени должна отвечать на внешние события в пределах заранее заданных временных ограничений.

\begin{definition}
\textbf{Дедлайн} задачи --- момент времени, к которому задача должна завершить выполнение.
\end{definition}

\begin{definition}
\textbf{Время отклика} --- интервал между поступлением события и завершением реакции системы на это событие.
\end{definition}

\begin{definition}
\textbf{Джиттер} --- отклонение фактического времени наступления события от ожидаемого или номинального времени.
\end{definition}

Например, если датчик должен опрашиваться каждые $10$ мс, но фактические интервалы равны $9.8$, $10.1$, $10.4$ мс, то разница между фактическими и идеальными моментами опроса называется джиттером.

\subsection{Классификация систем реального времени}

Системы реального времени обычно делят на три класса.

\subsubsection{Жёсткие системы реального времени}

\begin{definition}
\textbf{Жёсткая система реального времени} --- система, в которой нарушение дедлайна считается отказом системы.
\end{definition}

Примеры:
\begin{itemize}
    \item системы управления тормозами автомобиля;
    \item авиационные системы управления;
    \item медицинские имплантируемые устройства;
    \item системы аварийной защиты на производстве.
\end{itemize}

В жёстких системах важно гарантировать верхнюю оценку времени выполнения всех критических операций.

\subsubsection{Мягкие системы реального времени}

\begin{definition}
\textbf{Мягкая система реального времени} --- система, в которой нарушение дедлайна ухудшает качество работы, но не приводит к катастрофическому отказу.
\end{definition}

Примеры:
\begin{itemize}
    \item видеоконференции;
    \item потоковая передача аудио и видео;
    \item компьютерные игры;
    \item пользовательские интерфейсы.
\end{itemize}

Например, если кадр видео пришёл слишком поздно, его можно пропустить, но система продолжит работу.

\subsubsection{Твёрдые или firm real-time системы}

\begin{definition}
\textbf{Firm real-time система} --- система, в которой результат, полученный после дедлайна, бесполезен, но единичное нарушение дедлайна не обязательно является отказом всей системы.
\end{definition}

Пример: система обработки радиолокационных сигналов. Если отметка цели обработана слишком поздно, она может быть уже бесполезна, но последующие отметки всё ещё могут быть обработаны.

\subsection{Основные особенности программирования систем реального времени}

Программирование систем реального времени отличается от обычного прикладного программирования следующими особенностями.

\subsubsection{Временная корректность}

Обычная программа считается корректной, если она вычисляет правильный результат. Программа реального времени должна также сделать это вовремя.

Пусть задача $\tau_i$ имеет:
\begin{itemize}
    \item $C_i$ --- худшее время выполнения;
    \item $T_i$ --- период поступления;
    \item $D_i$ --- относительный дедлайн.
\end{itemize}

Тогда для каждого запуска задачи необходимо выполнение условия:
\[
    f_i \leq r_i + D_i,
\]
где:
\begin{itemize}
    \item $r_i$ --- момент поступления экземпляра задачи;
    \item $f_i$ --- момент завершения экземпляра задачи.
\end{itemize}

\subsubsection{Необходимость анализа худшего случая}

В системах реального времени важен не средний случай, а худший случай.

\begin{definition}
\textbf{WCET} от англ. \emph{Worst-Case Execution Time} --- верхняя оценка времени выполнения фрагмента программы в худшем случае.
\end{definition}

Если задача обычно выполняется за $2$ мс, но иногда может выполняться за $7$ мс, при анализе жёсткой системы реального времени необходимо использовать именно $7$ мс или доказанную верхнюю оценку.

\subsubsection{Детерминированность}

\begin{definition}
\textbf{Детерминированность} --- свойство системы, при котором её поведение предсказуемо и ограничено по времени.
\end{definition}

В системах реального времени нежелательны:
\begin{itemize}
    \item непредсказуемые задержки;
    \item динамическое выделение памяти в критических участках;
    \item неограниченные циклы;
    \item рекурсия без строгой глубины;
    \item блокирующие операции с неизвестным временем ожидания;
    \item сборка мусора без временных гарантий;
    \item страничная виртуальная память с непредсказуемыми page fault.
\end{itemize}

\subsubsection{Работа с внешними событиями}

Системы реального времени часто взаимодействуют с физическим миром:
\begin{itemize}
    \item датчиками;
    \item исполнительными механизмами;
    \item коммуникационными шинами;
    \item аппаратными таймерами;
    \item прерываниями.
\end{itemize}

Программа должна своевременно реагировать на события, например:
\begin{itemize}
    \item изменение уровня сигнала;
    \item поступление пакета по сети;
    \item срабатывание таймера;
    \item изменение показаний датчика;
    \item аварийное состояние.
\end{itemize}

\subsubsection{Ограниченность ресурсов}

Многие системы реального времени являются встроенными системами. Для них характерны:
\begin{itemize}
    \item малый объём оперативной памяти;
    \item отсутствие полноценной файловой системы;
    \item ограниченная частота процессора;
    \item отсутствие MMU;
    \item энергопотребление как важный критерий;
    \item необходимость работы без перезагрузки в течение длительного времени.
\end{itemize}

\subsubsection{Параллелизм и конкуренция}

Реальные системы обычно выполняют несколько задач:
\begin{itemize}
    \item опрос датчиков;
    \item фильтрацию данных;
    \item управление приводами;
    \item обмен по сети;
    \item диагностику;
    \item журналирование.
\end{itemize}

Поэтому возникают проблемы:
\begin{itemize}
    \item взаимного исключения;
    \item приоритетной инверсии;
    \item дедлоков;
    \item гонок данных;
    \item некорректного доступа к разделяемым ресурсам.
\end{itemize}

\subsection{Модель периодических задач}

Во многих теоретических моделях система реального времени описывается набором периодических задач:
\[
    \Gamma = \{\tau_1, \tau_2, \dots, \tau_n\}.
\]

Каждая задача $\tau_i$ задаётся параметрами:
\[
    \tau_i = (C_i, T_i, D_i),
\]
где:
\begin{itemize}
    \item $C_i$ --- худшее время выполнения одного экземпляра задачи;
    \item $T_i$ --- период;
    \item $D_i$ --- относительный дедлайн.
\end{itemize}

Если дедлайн равен периоду, то:
\[
    D_i = T_i.
\]

Коэффициент загрузки процессора задачей:
\[
    U_i = \frac{C_i}{T_i}.
\]

Суммарная загрузка процессора:
\[
    U = \sum_{i=1}^{n} \frac{C_i}{T_i}.
\]

Необходимое условие выполнимости на одном процессоре:
\[
    U \leq 1.
\]

Если $U > 1$, то даже при идеальном планировщике процессору не хватит времени для выполнения всех задач.

\subsection{Планирование задач реального времени}

\begin{definition}
\textbf{Планирование} --- выбор задачи, которая должна выполняться процессором в данный момент времени.
\end{definition}

Планировщики бывают:
\begin{itemize}
    \item вытесняющие;
    \item невытесняющие;
    \item со статическими приоритетами;
    \item с динамическими приоритетами;
    \item однопроцессорные;
    \item многопроцессорные.
\end{itemize}

\subsubsection{Вытесняющее и невытесняющее планирование}

При \textbf{вытесняющем планировании} задача может быть прервана более приоритетной задачей.

При \textbf{невытесняющем планировании} задача, начав выполнение, выполняется до завершения или добровольной блокировки.

Вытесняющее планирование даёт более быструю реакцию на критические события, но усложняет анализ и синхронизацию.

\subsection{Алгоритм Rate Monotonic Scheduling}

\begin{definition}
\textbf{Rate Monotonic Scheduling} --- алгоритм планирования периодических задач со статическими приоритетами, в котором чем меньше период задачи, тем выше её приоритет.
\end{definition}

Если:
\[
    T_i < T_j,
\]
то задача $\tau_i$ получает более высокий приоритет, чем $\tau_j$.

\subsubsection{Идея алгоритма}

Часто выполняющиеся задачи обычно требуют более быстрой реакции, поэтому им назначается больший приоритет.

\subsubsection{Шаги алгоритма Rate Monotonic Scheduling}

\begin{enumerate}
    \item Для каждой задачи $\tau_i$ определить период $T_i$.
    \item Отсортировать задачи по возрастанию периода.
    \item Задаче с наименьшим периодом назначить наивысший приоритет.
    \item Задаче со следующим по величине периодом назначить следующий приоритет.
    \item Во время выполнения всегда запускать готовую задачу с наибольшим приоритетом.
    \item Если появляется задача с более высоким приоритетом, текущая задача вытесняется.
\end{enumerate}

\subsubsection{Минимальный пример}

Пусть заданы две задачи:
\[
    \tau_1 = (C_1 = 1, T_1 = 4),
\]
\[
    \tau_2 = (C_2 = 2, T_2 = 6).
\]

Так как:
\[
    T_1 = 4 < T_2 = 6,
\]
задача $\tau_1$ получает более высокий приоритет.

Суммарная загрузка:
\[
    U = \frac{1}{4} + \frac{2}{6}
      = 0.25 + 0.333\ldots
      = 0.583\ldots
\]

Расписание на интервале $[0,12]$:

\[
\begin{array}{c|cccccccccccc}
t        & 0 & 1 & 2 & 3 & 4 & 5 & 6 & 7 & 8 & 9 & 10 & 11 \\
\hline
\text{CPU} & \tau_1 & \tau_2 & \tau_2 & - & \tau_1 & - & \tau_2 & \tau_2 & \tau_1 & - & - & -
\end{array}
\]

В моменты $0,4,8$ появляется задача $\tau_1$.
В моменты $0,6$ появляется задача $\tau_2$.

\subsection{Теорема о достаточном условии планируемости Rate Monotonic}

\begin{theorem}
Пусть имеется $n$ независимых периодических задач на одном процессоре. Для каждой задачи $\tau_i$ выполнено $D_i = T_i$. Задачи планируются вытесняющим алгоритмом Rate Monotonic. Если
\[
    \sum_{i=1}^{n} \frac{C_i}{T_i}
    \leq
    n(2^{1/n} - 1),
\]
то все задачи гарантированно укладываются в свои дедлайны.
\end{theorem}

\subsubsection{Пояснение}

Величина
\[
    U_{RM}(n) = n(2^{1/n} - 1)
\]
называется границей Лью и Лейланда для алгоритма Rate Monotonic.

При больших $n$:
\[
    \lim_{n \to \infty} n(2^{1/n} - 1) = \ln 2 \approx 0.693.
\]

То есть при большом числе задач достаточная гарантия планируемости по RM составляет примерно $69.3\%$ загрузки процессора.

\subsubsection{Доказательство}

Рассмотрим сначала две задачи:
\[
    \tau_1 = (C_1, T_1),
    \qquad
    \tau_2 = (C_2, T_2),
\]
где:
\[
    T_1 \leq T_2.
\]

По алгоритму Rate Monotonic задача $\tau_1$ имеет более высокий приоритет.

Для того чтобы задача $\tau_2$ успевала завершиться к своему дедлайну, за интервал длины $T_2$ должно хватить времени как на выполнение самой $\tau_2$, так и на все экземпляры более приоритетной задачи $\tau_1$.

В худшем случае обе задачи выпускаются одновременно в момент времени $0$. Тогда за интервал $[0,T_2]$ задача $\tau_1$ может появиться:
\[
    \left\lceil \frac{T_2}{T_1} \right\rceil
\]
раз.

Условие планируемости второй задачи:
\[
    C_2 + \left\lceil \frac{T_2}{T_1} \right\rceil C_1 \leq T_2.
\]

Это точное условие для двух задач, но оно содержит потолочную функцию. Лью и Лейланд получили более простую достаточную оценку.

Для общего случая $n$ задач рассматривается наихудшая конфигурация периодов, при которой задача с меньшим приоритетом испытывает максимальное влияние от задач с большими приоритетами.

Пусть задачи отсортированы по возрастанию периодов:
\[
    T_1 \leq T_2 \leq \dots \leq T_n.
\]

Для задачи $\tau_k$ помехи создают задачи:
\[
    \tau_1, \tau_2, \dots, \tau_{k-1}.
\]

В наихудшем случае критический момент наступает тогда, когда все задачи выпускаются одновременно.

Для периодических задач с дедлайнами, равными периодам, достаточным условием планируемости первых $k$ задач является:
\[
    \prod_{i=1}^{k} (U_i + 1) \leq 2,
\]
где:
\[
    U_i = \frac{C_i}{T_i}.
\]

По неравенству между средним арифметическим и средним геометрическим:
\[
    \prod_{i=1}^{k} (1 + U_i)
    \leq
    \left(
        1 + \frac{1}{k}\sum_{i=1}^{k} U_i
    \right)^k.
\]

Чтобы гарантировать:
\[
    \prod_{i=1}^{k} (1 + U_i) \leq 2,
\]
достаточно потребовать:
\[
    \left(
        1 + \frac{1}{k}\sum_{i=1}^{k} U_i
    \right)^k
    \leq 2.
\]

Извлечём корень степени $k$:
\[
    1 + \frac{1}{k}\sum_{i=1}^{k} U_i
    \leq
    2^{1/k}.
\]

Следовательно:
\[
    \frac{1}{k}\sum_{i=1}^{k} U_i
    \leq
    2^{1/k} - 1.
\]

Отсюда:
\[
    \sum_{i=1}^{k} U_i
    \leq
    k(2^{1/k} - 1).
\]

Для $k=n$ получаем:
\[
    \sum_{i=1}^{n} \frac{C_i}{T_i}
    \leq
    n(2^{1/n} - 1).
\]

Это и требовалось доказать.

\subsection{Алгоритм Earliest Deadline First}

\begin{definition}
\textbf{Earliest Deadline First} --- алгоритм планирования, в котором наивысший приоритет получает готовая задача с ближайшим абсолютным дедлайном.
\end{definition}

В отличие от Rate Monotonic, алгоритм EDF использует динамические приоритеты.

Абсолютный дедлайн экземпляра задачи:
\[
    d_{i,k} = r_{i,k} + D_i,
\]
где:
\begin{itemize}
    \item $r_{i,k}$ --- момент выпуска $k$-го экземпляра задачи $\tau_i$;
    \item $D_i$ --- относительный дедлайн.
\end{itemize}

\subsubsection{Шаги алгоритма EDF}

\begin{enumerate}
    \item При поступлении задачи вычислить её абсолютный дедлайн.
    \item Поместить задачу в очередь готовых задач.
    \item Отсортировать готовые задачи по возрастанию абсолютного дедлайна.
    \item Выполнять задачу с ближайшим дедлайном.
    \item Если поступила новая задача с более ранним дедлайном, вытеснить текущую задачу.
    \item После завершения задачи удалить её из очереди.
\end{enumerate}

\subsubsection{Минимальный пример}

Пусть имеются две задачи:
\[
    \tau_1 = (C_1 = 1, T_1 = 4, D_1 = 4),
\]
\[
    \tau_2 = (C_2 = 2, T_2 = 6, D_2 = 6).
\]

В момент $t=0$ появляются обе задачи:
\[
    d_1 = 0 + 4 = 4,
\]
\[
    d_2 = 0 + 6 = 6.
\]

Сначала выполняется $\tau_1$, так как её дедлайн раньше.

Расписание:
\[
\begin{array}{c|cccccccccccc}
t        & 0 & 1 & 2 & 3 & 4 & 5 & 6 & 7 & 8 & 9 & 10 & 11 \\
\hline
\text{CPU} & \tau_1 & \tau_2 & \tau_2 & - & \tau_1 & - & \tau_2 & \tau_2 & \tau_1 & - & - & -
\end{array}
\]

В этом примере расписание совпадает с RM, но в общем случае EDF может вести себя иначе.

\subsection{Теорема о планируемости EDF}

\begin{theorem}
Пусть имеется набор независимых периодических задач на одном процессоре. Для каждой задачи выполнено:
\[
    D_i = T_i.
\]
При вытесняющем планировании алгоритмом Earliest Deadline First набор задач планируем тогда и только тогда, когда:
\[
    \sum_{i=1}^{n} \frac{C_i}{T_i} \leq 1.
\]
\end{theorem}

\subsubsection{Доказательство}

Докажем необходимость.

Если:
\[
    \sum_{i=1}^{n} \frac{C_i}{T_i} > 1,
\]
то суммарная потребность задач в процессорном времени превышает доступное процессорное время. За достаточно большой интервал $[0,L]$ задачи потребуют более чем $L$ единиц процессорного времени. Один процессор не может выполнить больше $L$ единиц работы за интервал длины $L$. Следовательно, хотя бы один дедлайн будет нарушен. Необходимость доказана.

Докажем достаточность.

Пусть:
\[
    \sum_{i=1}^{n} \frac{C_i}{T_i} \leq 1.
\]

Предположим противное: алгоритм EDF нарушает некоторый дедлайн. Рассмотрим первый момент времени $t_d$, в который некоторый экземпляр задачи не завершился к своему дедлайну.

Пусть этот экземпляр был выпущен в момент $r$ и имел дедлайн $t_d$.

Рассмотрим интервал $[r,t_d]$. По правилу EDF в этом интервале процессор всегда отдаёт предпочтение задачам с дедлайнами не позже $t_d$. Если первый пропуск дедлайна произошёл в момент $t_d$, то до этого момента все задачи с более ранними дедлайнами выполнялись корректно.

Следовательно, пропуск дедлайна в $t_d$ возможен только в том случае, если суммарный объём работы задач с дедлайнами не позже $t_d$, выпущенных в рассматриваемом интервале, превышает длину интервала:
\[
    W(r,t_d) > t_d - r.
\]

Для периодических задач с $D_i = T_i$ максимальная работа, требуемая задачей $\tau_i$ на интервале длины $L$, не превосходит:
\[
    \frac{L}{T_i} C_i
\]
с учётом предельного случая на достаточно больших интервалах и синхронного выпуска задач.

Тогда суммарная работа ограничена:
\[
    W \leq L \sum_{i=1}^{n} \frac{C_i}{T_i}.
\]

Так как:
\[
    \sum_{i=1}^{n} \frac{C_i}{T_i} \leq 1,
\]
получаем:
\[
    W \leq L.
\]

Это противоречит необходимости $W > L$ для пропуска дедлайна.

Следовательно, при $U \leq 1$ алгоритм EDF строит корректное расписание. Достаточность доказана.

\subsection{Сравнение RM и EDF}

\[
\begin{array}{|c|c|c|}
\hline
\text{Свойство} & \text{Rate Monotonic} & \text{Earliest Deadline First} \\
\hline
\text{Тип приоритетов} & \text{Статические} & \text{Динамические} \\
\hline
\text{Критерий приоритета} & \text{Меньший период} & \text{Ближайший дедлайн} \\
\hline
\text{Максимальная загрузка} & \approx 69.3\% \text{ гарантированно} & 100\% \\
\hline
\text{Простота реализации} & \text{Высокая} & \text{Ниже} \\
\hline
\text{Предсказуемость} & \text{Высокая} & \text{Требует анализа} \\
\hline
\text{Использование в RTOS} & \text{Очень часто} & \text{Реже} \\
\hline
\end{array}
\]

\subsection{Приоритетная инверсия}

\begin{definition}
\textbf{Приоритетная инверсия} --- ситуация, в которой высокоприоритетная задача вынуждена ждать низкоприоритетную задачу из-за захваченного разделяемого ресурса.
\end{definition}

\subsubsection{Пример приоритетной инверсии}

Пусть есть три задачи:
\[
    H \quad \text{--- высокий приоритет},
\]
\[
    M \quad \text{--- средний приоритет},
\]
\[
    L \quad \text{--- низкий приоритет}.
\]

Пусть задача $L$ захватила мьютекс $R$.

Затем появляется задача $H$ и тоже хочет захватить $R$, но ресурс занят. Поэтому $H$ блокируется.

После этого начинает выполняться задача $M$, так как она имеет приоритет выше, чем $L$, но ниже, чем $H$.

В результате:
\[
    H \text{ ждёт } L,
\]
но $L$ не может продолжить выполнение, потому что её вытесняет $M$.

Получается, что среднеприоритетная задача $M$ фактически задерживает высокоприоритетную задачу $H$.

\subsubsection{Иллюстрация}

\[
\begin{array}{c|c|l}
\text{Время} & \text{Событие} & \text{Комментарий} \\
\hline
t_0 & L \text{ захватывает } R & L \text{ выполняется} \\
t_1 & H \text{ требует } R & H \text{ блокируется} \\
t_2 & M \text{ становится готовой} & M \text{ вытесняет } L \\
t_3 & M \text{ выполняется} & H \text{ всё ещё ждёт} \\
t_4 & L \text{ продолжает работу} & L \text{ освобождает } R \\
t_5 & H \text{ получает } R & H \text{ продолжает выполнение}
\end{array}
\]

\subsection{Протокол наследования приоритета}

\begin{definition}
\textbf{Протокол наследования приоритета} --- механизм, при котором низкоприоритетная задача, удерживающая ресурс, временно наследует приоритет высокоприоритетной задачи, ожидающей этот ресурс.
\end{definition}

\subsubsection{Шаги протокола}

\begin{enumerate}
    \item Задача $L$ захватывает ресурс $R$.
    \item Задача $H$ с более высоким приоритетом пытается захватить $R$.
    \item Так как $R$ занят, задача $H$ блокируется.
    \item Задача $L$ временно получает приоритет задачи $H$.
    \item Задача $L$ завершает критическую секцию без вытеснения среднеприоритетными задачами.
    \item Задача $L$ освобождает ресурс $R$.
    \item Задача $H$ разблокируется.
    \item Задача $L$ возвращается к своему исходному приоритету.
\end{enumerate}

\subsection{Протокол потолка приоритетов}

\begin{definition}
\textbf{Протокол потолка приоритетов} --- механизм синхронизации, при котором каждому ресурсу заранее назначается потолочный приоритет, равный максимальному приоритету среди задач, которые могут использовать этот ресурс.
\end{definition}

Если задача захватывает ресурс, её приоритет может быть поднят до потолка этого ресурса.

Преимущества:
\begin{itemize}
    \item ограничение времени блокировки;
    \item предотвращение некоторых дедлоков;
    \item упрощение анализа времени отклика.
\end{itemize}

\subsection{Анализ времени отклика}

Для задач со статическими приоритетами часто используется анализ времени отклика.

Пусть задачи отсортированы по убыванию приоритета. Для задачи $\tau_i$ время отклика $R_i$ можно найти итерационно:
\[
    R_i^{(0)} = C_i,
\]
\[
    R_i^{(k+1)}
    =
    C_i + B_i +
    \sum_{\tau_j \in hp(i)}
    \left\lceil
        \frac{R_i^{(k)}}{T_j}
    \right\rceil C_j,
\]
где:
\begin{itemize}
    \item $hp(i)$ --- множество задач с приоритетом выше, чем у $\tau_i$;
    \item $B_i$ --- максимальное время блокировки задачи $\tau_i$ низкоприоритетными задачами;
    \item $C_i$ --- собственное время выполнения задачи;
    \item сумма описывает помехи от более приоритетных задач.
\end{itemize}

Задача планируема, если:
\[
    R_i \leq D_i.
\]

\subsubsection{Алгоритм вычисления времени отклика}

\begin{enumerate}
    \item Установить начальное приближение:
    \[
        R_i^{(0)} = C_i + B_i.
    \]
    \item На каждой итерации вычислять:
    \[
        R_i^{(k+1)}
        =
        C_i + B_i +
        \sum_{\tau_j \in hp(i)}
        \left\lceil
            \frac{R_i^{(k)}}{T_j}
        \right\rceil C_j.
    \]
    \item Если:
    \[
        R_i^{(k+1)} = R_i^{(k)},
    \]
    то найдено фиксированное значение времени отклика.
    \item Если:
    \[
        R_i^{(k+1)} > D_i,
    \]
    то задача не удовлетворяет дедлайну.
    \item Повторить проверку для всех задач.
\end{enumerate}

\subsubsection{Минимальный пример}

Пусть:
\[
    \tau_1 = (C_1=1, T_1=4, D_1=4),
\]
\[
    \tau_2 = (C_2=2, T_2=10, D_2=10).
\]

Задача $\tau_1$ имеет более высокий приоритет.

Для $\tau_1$:
\[
    R_1 = C_1 = 1 \leq 4.
\]

Для $\tau_2$:
\[
    R_2^{(0)} = C_2 = 2.
\]

Первая итерация:
\[
    R_2^{(1)}
    =
    2 + \left\lceil \frac{2}{4} \right\rceil \cdot 1
    =
    2 + 1 = 3.
\]

Вторая итерация:
\[
    R_2^{(2)}
    =
    2 + \left\lceil \frac{3}{4} \right\rceil \cdot 1
    =
    3.
\]

Получили неподвижную точку:
\[
    R_2 = 3.
\]

Так как:
\[
    R_2 = 3 \leq D_2 = 10,
\]
задача $\tau_2$ укладывается в дедлайн.

\subsection{Прерывания в системах реального времени}

\begin{definition}
\textbf{Прерывание} --- аппаратный или программный сигнал, временно приостанавливающий текущее выполнение процессора для обработки срочного события.
\end{definition}

Прерывания позволяют системе быстро реагировать на внешние события, но требуют осторожного проектирования.

\subsubsection{Правила обработки прерываний}

В обработчиках прерываний желательно:
\begin{itemize}
    \item выполнять только минимально необходимую работу;
    \item не использовать длительные циклы;
    \item не выполнять блокирующие операции;
    \item не вызывать функции с непредсказуемым временем выполнения;
    \item передавать основную обработку обычной задаче;
    \item минимизировать время с запрещёнными прерываниями.
\end{itemize}

Типовой подход:
\begin{enumerate}
    \item Прерывание фиксирует факт события.
    \item Обработчик быстро считывает минимальные данные.
    \item Обработчик уведомляет задачу через очередь, семафор или флаг.
    \item Основная обработка выполняется в задаче.
\end{enumerate}

\subsubsection{Минимальный пример на псевдокоде}

\begin{verbatim}
interrupt_handler UART_RX_ISR:
    byte = UART.read()
    queue_push(rx_queue, byte)
    signal(rx_task)

task rx_task:
    while true:
        wait_signal()
        while not queue_empty(rx_queue):
            byte = queue_pop(rx_queue)
            process(byte)
\end{verbatim}

Обработчик прерывания только считывает байт и помещает его в очередь. Более длительная обработка выполняется в задаче.

\subsection{Синхронизация задач}

В системах реального времени используются следующие средства синхронизации:
\begin{itemize}
    \item мьютексы;
    \item бинарные семафоры;
    \item счётные семафоры;
    \item очереди сообщений;
    \item event flags;
    \item spinlock;
    \item lock-free структуры данных.
\end{itemize}

\subsubsection{Мьютекс}

\begin{definition}
\textbf{Мьютекс} --- объект взаимного исключения, позволяющий только одной задаче одновременно находиться в критической секции.
\end{definition}

Пример:
\begin{verbatim}
lock(mutex)
shared_counter = shared_counter + 1
unlock(mutex)
\end{verbatim}

В системах реального времени мьютекс должен иметь ограниченное время ожидания и желательно поддерживать протокол наследования приоритета.

\subsubsection{Семафор}

\begin{definition}
\textbf{Семафор} --- синхронизационный объект, содержащий счётчик доступных разрешений.
\end{definition}

Бинарный семафор часто используется для уведомления задачи о событии.

Пример:
\begin{verbatim}
interrupt_handler TIMER_ISR:
    give(semaphore)

task control_task:
    while true:
        take(semaphore)
        control_step()
\end{verbatim}

\subsection{Память в системах реального времени}

Обычные методы управления памятью могут быть опасны в системах реального времени.

Проблемы динамического выделения памяти:
\begin{itemize}
    \item непредсказуемое время выполнения \texttt{malloc};
    \item фрагментация памяти;
    \item невозможность выделить память в критический момент;
    \item скрытые блокировки внутри менеджера памяти.
\end{itemize}

Поэтому часто применяют:
\begin{itemize}
    \item статическое выделение памяти;
    \item memory pool;
    \item ring buffer;
    \item заранее выделенные очереди;
    \item запрет динамической памяти после инициализации.
\end{itemize}

\subsubsection{Кольцевой буфер}

\begin{definition}
\textbf{Кольцевой буфер} --- массив фиксированного размера, в котором индексы чтения и записи циклически переходят к началу массива.
\end{definition}

Пусть размер буфера равен $N$.

Индекс записи:
\[
    head = (head + 1) \bmod N.
\]

Индекс чтения:
\[
    tail = (tail + 1) \bmod N.
\]

\subsubsection{Алгоритм записи в кольцевой буфер}

\begin{enumerate}
    \item Вычислить следующий индекс:
    \[
        next = (head + 1) \bmod N.
    \]
    \item Если:
    \[
        next = tail,
    \]
    то буфер полон.
    \item Иначе записать элемент:
    \[
        buffer[head] = x.
    \]
    \item Обновить индекс:
    \[
        head = next.
    \]
\end{enumerate}

\subsubsection{Алгоритм чтения из кольцевого буфера}

\begin{enumerate}
    \item Если:
    \[
        head = tail,
    \]
    то буфер пуст.
    \item Иначе прочитать:
    \[
        x = buffer[tail].
    \]
    \item Обновить индекс:
    \[
        tail = (tail + 1) \bmod N.
    \]
\end{enumerate}

\subsection{Типичная архитектура программы реального времени}

Типовая программа реального времени может иметь следующую структуру:

\[
\begin{array}{c}
\text{Аппаратные устройства} \\
\downarrow \\
\text{Прерывания и драйверы} \\
\downarrow \\
\text{RTOS: планировщик, таймеры, очереди, семафоры} \\
\downarrow \\
\text{Прикладные задачи} \\
\downarrow \\
\text{Алгоритмы управления и обработки данных}
\end{array}
\]

\subsubsection{Пример набора задач}

\[
\begin{array}{|c|c|c|c|}
\hline
\text{Задача} & C_i & T_i & \text{Назначение} \\
\hline
\tau_1 & 1 \text{ мс} & 5 \text{ мс} & \text{опрос датчика} \\
\hline
\tau_2 & 2 \text{ мс} & 10 \text{ мс} & \text{управление приводом} \\
\hline
\tau_3 & 5 \text{ мс} & 50 \text{ мс} & \text{передача данных} \\
\hline
\end{array}
\]

Суммарная загрузка:
\[
    U =
    \frac{1}{5}
    +
    \frac{2}{10}
    +
    \frac{5}{50}
    =
    0.2 + 0.2 + 0.1
    =
    0.5.
\]

Для $n=3$ граница RM:
\[
    U_{RM}(3)
    =
    3(2^{1/3} - 1)
    \approx
    0.779.
\]

Так как:
\[
    0.5 \leq 0.779,
\]
набор задач гарантированно планируем алгоритмом RM.

\subsection{Реактивный стиль программирования}

Системы реального времени часто являются реактивными.

\begin{definition}
\textbf{Реактивная система} --- система, которая постоянно взаимодействует с окружающей средой и реагирует на входные события.
\end{definition}

Реактивная программа обычно имеет вид:
\begin{verbatim}
initialize()

while true:
    event = wait_event()
    handle(event)
\end{verbatim}

Для простых встроенных систем часто применяется архитектура \emph{superloop}:

\begin{verbatim}
initialize()

while true:
    read_sensors()
    compute_control()
    update_actuators()
    communicate()
\end{verbatim}

Преимущества superloop:
\begin{itemize}
    \item простота;
    \item отсутствие сложного планировщика;
    \item малые накладные расходы.
\end{itemize}

Недостатки:
\begin{itemize}
    \item сложно гарантировать время реакции при росте сложности;
    \item одна долгая операция задерживает остальные;
    \item трудно добавлять задачи с разными периодами;
    \item неудобно использовать блокирующий ввод-вывод.
\end{itemize}

\subsection{Конечные автоматы}

Для описания поведения систем реального времени часто применяются конечные автоматы.

\begin{definition}
\textbf{Конечный автомат} --- модель вычислений, состоящая из конечного множества состояний, событий и переходов между состояниями.
\end{definition}

Формально автомат можно задать пятёркой:
\[
    A = (S, E, \delta, s_0, F),
\]
где:
\begin{itemize}
    \item $S$ --- множество состояний;
    \item $E$ --- множество событий;
    \item $\delta: S \times E \to S$ --- функция переходов;
    \item $s_0$ --- начальное состояние;
    \item $F$ --- множество завершающих или особых состояний.
\end{itemize}

\subsubsection{Пример автомата}

Рассмотрим простой контроллер двери.

Состояния:
\[
    S = \{\text{Closed}, \text{Opening}, \text{Open}, \text{Closing}\}.
\]

События:
\[
    E = \{\text{button}, \text{opened}, \text{closed}, \text{timeout}\}.
\]

Таблица переходов:

\[
\begin{array}{|c|c|c|}
\hline
\text{Текущее состояние} & \text{Событие} & \text{Новое состояние} \\
\hline
\text{Closed} & \text{button} & \text{Opening} \\
\hline
\text{Opening} & \text{opened} & \text{Open} \\
\hline
\text{Open} & \text{timeout} & \text{Closing} \\
\hline
\text{Closing} & \text{closed} & \text{Closed} \\
\hline
\end{array}
\]

Псевдокод:

\begin{verbatim}
state = Closed

on_event(event):
    if state == Closed and event == button:
        motor_open()
        state = Opening

    elif state == Opening and event == opened:
        motor_stop()
        state = Open
        start_timer()

    elif state == Open and event == timeout:
        motor_close()
        state = Closing

    elif state == Closing and event == closed:
        motor_stop()
        state = Closed
\end{verbatim}

Конечные автоматы удобны тем, что делают поведение системы явным и проверяемым.

\subsection{Таймеры и периодические задачи}

В системах реального времени широко используются:
\begin{itemize}
    \item аппаратные таймеры;
    \item программные таймеры;
    \item системный тик;
    \item tickless-режимы;
    \item watchdog-таймеры.
\end{itemize}

\begin{definition}
\textbf{Watchdog-таймер} --- аппаратный или программный таймер, который перезапускает систему или переводит её в безопасное состояние, если программа своевременно не подтверждает свою работоспособность.
\end{definition}

Типовая схема:

\begin{verbatim}
while true:
    do_critical_work()
    if system_is_healthy():
        watchdog_kick()
\end{verbatim}

Если программа зависла и не вызвала \texttt{watchdog\_kick()}, watchdog инициирует восстановление.

\subsection{Проблема блокировок}

Блокирующие операции опасны, если время ожидания не ограничено.

Примеры нежелательных операций в критических задачах:
\begin{itemize}
    \item ожидание сетевого ответа без тайм-аута;
    \item ожидание мьютекса без верхней оценки;
    \item запись во флеш-память с непредсказуемой длительностью;
    \item ввод-вывод через медленную шину;
    \item ожидание пользовательского ввода.
\end{itemize}

В системах реального времени следует использовать:
\begin{itemize}
    \item неблокирующие вызовы;
    \item тайм-ауты;
    \item разделение критической и некритической обработки;
    \item приоритетные очереди;
    \item асинхронную обработку событий.
\end{itemize}

\subsection{Ошибки, характерные для систем реального времени}

К типичным ошибкам относятся:
\begin{itemize}
    \item отсутствие анализа WCET;
    \item неправильный выбор приоритетов;
    \item слишком длинные обработчики прерываний;
    \item неограниченное ожидание мьютекса;
    \item приоритетная инверсия;
    \item использование динамической памяти в критическом пути;
    \item неограниченный рост очередей;
    \item отсутствие обработки переполнения буферов;
    \item недооценка джиттера;
    \item отсутствие watchdog;
    \item попытка проверять систему только тестированием без аналитических гарантий.
\end{itemize}

\subsection{Тестирование и верификация}

Для систем реального времени применяются:
\begin{itemize}
    \item модульное тестирование;
    \item интеграционное тестирование;
    \item нагрузочное тестирование;
    \item стресс-тестирование;
    \item измерение времени выполнения;
    \item статический анализ;
    \item формальная верификация;
    \item моделирование расписаний;
    \item трассировка выполнения.
\end{itemize}

Важно различать:
\[
    \text{измеренное время выполнения}
    \neq
    \text{доказанное худшее время выполнения}.
\]

Измерения показывают только то, что произошло в тестовых сценариях. Для жёстких систем реального времени требуется верхняя оценка худшего случая.

\subsection{Основные принципы проектирования программ реального времени}

\begin{enumerate}
    \item Сначала определить временные требования.
    \item Разделить систему на задачи с известными периодами и дедлайнами.
    \item Оценить или доказать WCET каждой задачи.
    \item Выбрать алгоритм планирования.
    \item Проверить планируемость.
    \item Минимизировать критические секции.
    \item Избегать неограниченных блокировок.
    \item Использовать статическое или заранее ограниченное выделение памяти.
    \item Делать обработчики прерываний короткими.
    \item Использовать протоколы борьбы с приоритетной инверсией.
    \item Предусматривать безопасное состояние при отказе.
    \item Документировать все временные предположения.
\end{enumerate}

\subsection{Краткий итог}

Программирование систем реального времени ориентировано на обеспечение гарантированных временных характеристик. Главная особенность таких систем состоит в том, что правильный результат, полученный слишком поздно, может быть бесполезен или опасен.

Ключевые понятия:
\begin{itemize}
    \item дедлайн;
    \item время отклика;
    \item джиттер;
    \item WCET;
    \item планируемость;
    \item приоритеты;
    \item прерывания;
    \item приоритетная инверсия;
    \item протоколы синхронизации;
    \item детерминированность.
\end{itemize}

Основные алгоритмы планирования:
\begin{itemize}
    \item Rate Monotonic --- статические приоритеты по периодам;
    \item Earliest Deadline First --- динамические приоритеты по ближайшим дедлайнам.
\end{itemize}

Для Rate Monotonic достаточное условие планируемости:
\[
    \sum_{i=1}^{n} \frac{C_i}{T_i}
    \leq
    n(2^{1/n} - 1).
\]

Для Earliest Deadline First при $D_i=T_i$ условие планируемости:
\[
    \sum_{i=1}^{n} \frac{C_i}{T_i}
    \leq 1.
\]

Главная инженерная цель состоит в том, чтобы система была не просто быстрой, а предсказуемой.

\end{document}
